\section{Second variation on background classes}
\label{sec:gaud-volume-hessian}

Differentiating the smooth equation twice along an affine family of backgrounds computes the second derivative of the root volume. After integration against the adjoint probability of the solution form, the second derivative of the potential disappears, and what remains is the integrated squared norm of the linearized solution form. Consequently the Hessian of $-\log\lambda_\mu$ is positive definite once smooth $\ddc$ directions are factored out. On the cone of closed K\"ahler classes, where $\lambda_\mu(\theta)=(\int_X\theta^n)^{1/n}$, it is $1/n$ times the Hessian of $-\log\int_X\theta^n$.

Fix $\theta\in\mathscr P$ and write $\Omega=\Omega_\theta$, $u=u_\theta$ and $\lambda=\lambda_\mu(\theta)$. Let $\nu_\Omega$ be the adjoint probability of $\Omega$ from \cref{prop:adjoint-gauduchon}, that is, the smooth positive probability measure, given by the Gauduchon multiplier of $\Omega$, which satisfies $\int_X\operatorname{tr}_\Omega(\ddc v)\,\mathrm{d}\nu_\Omega=0$ for every smooth real $v$. In dimension one this probability is $\Omega/\int_X\Omega=\mu$.

Let $\theta_t=\theta+t\chi$, where $t$ is real and $\chi$ is a smooth real $(1,1)$-form on $X$. For sufficiently small $|t|$, the background $\theta_t$ stays in $\mathscr P$. We normalize the solution potentials $u_t$ by one fixed condition on their mean, and dots denote derivatives of $u_t$ and of $\lambda_\mu(\theta_t)$ at $t=0$. Define the Hermitian endomorphism
\[
 A_\chi=\Omega^{-1}(\chi+\ddc\dot u)
\]
of $T^{1,0}X$, where $\Omega^{-1}$ means raising an index with the positive form $\Omega$; traces of powers of $A_\chi$ are independent of coordinates. We write $I$ for the identity endomorphism of $T^{1,0}X$ and $\|\cdot\|_\Omega$ for the Hilbert--Schmidt norm on its endomorphisms induced by $\Omega$.

\Needspace{16\baselineskip}
\begin{proposition}\label{prop:gaud-volume-hessian}\index{Gauduchon variational volume!Hessian}
Along this affine variation,
\begin{align}
 \frac{\ddot\lambda}{\lambda}
   &=-\frac1n\int_X
      \left\|A_\chi-\frac{\dot\lambda}{\lambda}I\right\|_\Omega^2
                                      \,\mathrm{d}\nu_\Omega,\label{eq:gaud-root-hessian}\\
 -\frac{\mathrm{d}^2}{\mathrm{d}t^2}\log\lambda_\mu(\theta_t)
   &=\frac1n\int_X\operatorname{tr}(A_\chi^2)\,\mathrm{d}\nu_\Omega.
                                      \label{eq:gaud-log-hessian}
\end{align}
The first expression vanishes exactly when
$\chi=r\theta+\ddc f$ for a real constant $r$ and a smooth real $f$.
The second vanishes exactly when $\chi=\ddc f$ for a smooth real $f$.
\end{proposition}
\begin{proof}
Smooth dependence on $t$ was proved in \cref{prop:gaud-volume-first-variation}. We differentiate the logarithmic form $\log(\Omega_t^n/\mu)=n\log\lambda_\mu(\theta_t)$ of the equation. Since $\mu$ is fixed,
\[
 \operatorname{tr}A_\chi=n\frac{\dot\lambda}{\lambda}.
\]
In particular this trace is constant on $X$. Differentiating once more along the affine path $\theta_t$ gives
\[
 \operatorname{tr}_\Omega(\ddc\ddot u)
       -\operatorname{tr}(A_\chi^2)
   =n\left(\frac{\ddot\lambda}{\lambda}
                  -\frac{\dot\lambda^2}{\lambda^2}\right).
\]
Integration against the probability $\nu_\Omega$, which is fixed since everything is evaluated at $t=0$, removes the first term. This proves \eqref{eq:gaud-log-hessian}.

Put $r=\dot\lambda/\lambda$. Since $A_\chi$ is Hermitian and $\operatorname{tr}A_\chi=nr$,
\[
 \|A_\chi-rI\|_\Omega^2
      =\operatorname{tr}(A_\chi^2)-nr^2.
\]
Substituting this into \eqref{eq:gaud-log-hessian} gives \eqref{eq:gaud-root-hessian}. Since the adjoint measure has a strictly positive smooth density, the right side of \eqref{eq:gaud-root-hessian} vanishes exactly when $A_\chi=rI$ everywhere, that is, when
\[
 \chi+\ddc\dot u=r(\theta+\ddc u).
\]
This is the stated relation with $f=ru-\dot u$. Conversely, if $\chi=r\theta+\ddc f$, then $\theta_t$ differs from $(1+rt)\theta$ by a smooth $\ddc$ term for small $t$. Homogeneity gives $\lambda_\mu(\theta_t)=(1+rt)\lambda$, whose second derivative is zero.

The integrand in \eqref{eq:gaud-log-hessian} is the squared norm of the Hermitian endomorphism $A_\chi$. Its integral therefore vanishes exactly when $A_\chi=0$, that is, when $\chi=-\ddc\dot u$. Conversely, $\lambda_\mu(\theta+t\ddc f)=\lambda$ for all $t$. This proves both equality statements.
\end{proof}

The linearized solution and the derivative of the constant depend linearly on $\chi$. For smooth real $(1,1)$-forms $\chi_1,\chi_2$ on $X$, polarization therefore gives a symmetric bilinear form
\begin{equation}\label{eq:gaud-background-quadratic-form}
 \mathfrak g_\theta(\chi_1,\chi_2)
    =-D^2\log\lambda_\mu|_\theta(\chi_1,\chi_2)
    =\frac1n\int_X\operatorname{tr}(A_{\chi_1}A_{\chi_2})\,\mathrm{d}\nu_\Omega.
\end{equation}
We call the classes of smooth real $(1,1)$-forms modulo $\ddc C^\infty(X,\mathbb R)$ \emph{background classes}. By \cref{prop:gaud-volume-hessian}, the form $\mathfrak g_\theta$ descends to background classes of directions and is positive definite there. A change of representative of $\theta$ leaves its solution form $\Omega$ unchanged, and a change of representative of a direction is absorbed by the derivative of the solution potential. Thus $\mathfrak g$ is intrinsically defined on the background classes in $\mathscr P$. On every finite-dimensional affine family whose directions are linearly independent modulo smooth $\ddc$ terms, it is a smooth Riemannian metric.

These variations change the background class. Directions obtained by varying a potential inside one fixed class are exactly the directions annihilated by \eqref{eq:gaud-background-quadratic-form}. Thus this construction concerns a different space from the potential-space geometry of \cref{chap:energy}.

\Needspace{13\baselineskip}
\section{Recovering a positive class from its Gauduchon derivative}
\label{sec:gaud-volume-gradient-recovery}

By \cref{prop:gaud-volume-first-variation}, the derivative of $\lambda_\mu$ at $\theta\in\mathscr P$ is represented by a minimizing Gauduchon metric, which is determined up to a constant factor. The solution form is conformal to this metric and has a prescribed volume form, so the value and the derivative of $\lambda_\mu$ together reconstruct the solution form, and hence the background class. Fix $\theta\in\mathscr P$ and a minimizing metric $\gamma_\theta$.

For $n\geq2$ define the smooth positive $(n-1,n-1)$-form
\[
 \Xi_\theta=\frac{\gamma_\theta^{n-1}}{D_\mu(\gamma_\theta)}.
\]
Since numerator and denominator are both homogeneous of degree $n-1$ in $\gamma_\theta$, this form does not depend on the choice of the minimizing metric. It is $\ddc$-closed, and by \cref{prop:gaud-volume-first-variation} it represents the derivative in every smooth real $(1,1)$-direction $\chi$:
\[
 (\mathrm{d}\lambda_\mu)_\theta(\chi)=\int_X\chi\wedge\Xi_\theta.
\]

\begin{proposition}\label{prop:gaud-derivative-recovery}
Let $n\geq2$ and $\theta\in\mathscr P$. The value $\lambda_\mu(\theta)$ and the derivative of $\lambda_\mu$ at $\theta$ on all smooth real $(1,1)$-forms together determine the background class of $\theta$, and the derivative alone determines this class up to multiplication by a positive constant. More precisely, if $\widehat\gamma$ is the positive $(1,1)$-form with $\widehat\gamma^{n-1}=\Xi_\theta$, then
\[
 D_\mu(\widehat\gamma)=1,
 \qquad
 \Omega_\theta=
   \lambda_\mu(\theta)
       \left(\frac{\mu}{\widehat\gamma^n}\right)^{1/n}
       \widehat\gamma.
\]
Moreover, assigning to a background class in $\mathscr P$ with $\lambda_\mu=1$ the ray of its minimizing Gauduchon metrics is a bijection onto the set of rays of Gauduchon metrics.
\end{proposition}
\begin{proof}
The values of the derivative on all smooth $(1,1)$-forms determine the smooth form $\Xi_\theta$. Its positive root exists and is unique by the positive-root construction preceding the proof of \cref{thm:lamari}. If $D=D_\mu(\gamma_\theta)$, this root is $\widehat\gamma=D^{-1/(n-1)}\gamma_\theta$, and the homogeneity of $D_\mu$ gives $D_\mu(\widehat\gamma)=1$.

By \cref{thm:gaud-volume-dual-formula}, the solution form $\Omega_\theta$ is conformal to $\widehat\gamma$, and its volume form is $\lambda_\mu(\theta)^n\mu$. These two requirements determine the conformal factor pointwise and give the formula for $\Omega_\theta$. Thus the value and the derivative recover $\Omega_\theta$, a representative of the class of $\theta$. Without the value, the same formula recovers $\Omega_\theta$ up to a positive constant factor. Conversely, multiplying the class by a positive constant multiplies $\Omega_\theta$ by the same constant and leaves its minimizing Gauduchon ray, and hence $\Xi_\theta$, unchanged; so the derivative cannot determine more than the ray of the class.

For the surjectivity in the last assertion, take any Gauduchon metric $\gamma$ and set
\[
 \Omega_\gamma=\left(\frac\mu{\gamma^n}\right)^{1/n}\gamma.
\]
This Hermitian form is unchanged by constant rescaling of $\gamma$ and satisfies $\Omega_\gamma^n=\mu$. The zero potential therefore solves the equation for the background $\Omega_\gamma$ with constant one, so its class has $\lambda_\mu=1$, and $\gamma$ is a minimizing Gauduchon metric by \cref{thm:gaud-volume-dual-formula}. Injectivity follows from the recovery formula: two classes with $\lambda_\mu=1$ and the same Gauduchon ray have the same solution form. This proves the bijection.
\end{proof}

In dimension one, a background class is determined by its integral, by the Green operator of the curve. The derivative is then the fixed functional $\chi\mapsto\int_X\chi$, and the value alone determines the class.
